\chapter{Felhasználói dokumentáció}
\label{ch:user}
\todo{Felhasználói dokumentáció}


Procedurális növényzet generálására több létező megoldás is található attól függően, hogy milyen céljaink, illetve kritériumaink vannak. Szempont lehet, hogy a vegetáció mennyire ítélhető realisztikusnak, mekkora területet fed le, milyen környezeti befolyásoló tényezők hatnak a növényekre és azok terjedésére. Ezen követelmények rendkívül fontosak lehetnek nagy, nyíltterű videójátékoknál, szimulációs szoftvereknél, illetve egyéb olyan eszközöknél, ahol egy-egy létező tájat szeretnék életszerűen lemodellezni. Minél valósabb eredményeket szeretnék elérni, annál több információ szükséges az adott módszerhez. Ezért is érdemes lehet különböző multidiszciplináris módszerekkel újabb forrásokat alkalmazni, és az ezeket körülvevő ökoszisztémára támaszkodni. 

Természetesen adódik így, hogy az előbb leírt problémára érdemes olyan eszközöket kézbe venni, amelyek alaból egy-egy terület helyezkedéséről, növényvilágáról és felbontásáról adnak információt tömören, de egyértelműen. A tájfutótérképek pont ilyen elvárásokat betartva vannak megtervezve. Ezek a fajta térképek egy-egy tetszőleges tájat írnak le a tájfutó versenyzők számára, jelölve a területen megtalálható természetes határokat, érdekességeket, és ember-által készített struktúrákat. Rendeltetésük miatt meg van határozva, hogy mekkora precizitással számolhatnak a felhasználóik. Továbbá a jelölések sztandardizálva van, így meg adva, hogy egy szimbólum mikor és mikor nem használható az egyes esetekben.\todo{Ezt a mondatot gondoljuk át} Ez egészen értelmezhető és kellően részletes keretrendszert nyújt számunkra ahhoz, hogy a térkép által adott információból különböző növényekkel tarkított területeket tudjuk meghatározni és egyedeit kellő képpen leírni.

Az, hogy milyen szinten szeretnénk a növényzetünket leírni teljesen mértékben applikációtól függ. Témakörömhöz így pár fő szempont lett elsősorban meghatározva. Ezek elsősorban a nagyobb cserjék, illetve fásszárú növények elhelyezkedése, az általuk elfoglalt tér meghatározása volt.





Vegetáció generálás komplexitása -> multidiszciplináris tudás -> térképeszetből merítünk ihletet



Megoldott probléma rövid leírására




Gépigény

Min. 8 Gb RAM



Telepítés